\itemhead{1.2 思考 · 未编号题组（教材p.19）}{XJ-B1-C01-TH-P019-Y1033}
\label{XJ-B1-C01-TH-P019-Y1033}
{\footnotesize 湘教版教材p.19 / PDF第26页 / 1.2 思考 / 未编号题组}
\field{原题}{下列语句是命题吗？将（2）（3）分别与（1）比较，你发现了什么？
\subq{1}$x>0$；
\subq{2}对每一个实数 $x$，有 $x>0$；
\subq{3}有一个实数 $x$，使 $x>0$。}
\field{参考答案}{（1）不是命题；（2）（3）都是命题，前者是假命题，后者是真命题。量词限定了变量的范围及断言方式。}
{\small \field{关联知识}{L03 全称、存在量词及真假的判定（教材帮p.38、39 / PDF47、48）}
\field{方法与易错点}{先确认量词、作用对象及论域，再采用任意性证明或实例反例。 易错：用几个满足的数证明全称命题；把“存在”理解为“唯一存在”。}
\field{答案依据}{答案册未提供本栏目的独立对应答案；依据原题、教材概念补写参考回应。开放性总结不设唯一数值答案。}}
